\section{Conclusiones}

La aplicación de la matriz insumo-producto de Leontief a la economía de La Guajira ha permitido ilustrar el uso y las bondades de esta herramienta para analizar las interacciones entre sectores económicos. A partir de los tres ejercicios desarrollados, se pueden extraer las siguientes conclusiones:

\begin{enumerate}
    \item \textbf{La matriz de Leontief es una herramienta poderosa para el análisis económico regional.} Permite cuantificar las interacciones entre sectores y evaluar el impacto de cambios en la demanda final sobre la producción total requerida.

    \item \textbf{Los multiplicadores de producción revelan la importancia estratégica de cada sector.} En el caso de La Guajira, el sector Tecnología presenta el mayor multiplicador (1.67), seguido de Agricultura (1.62) y Servicios (1.60), lo que indica que un aumento en la demanda de Tecnología genera el mayor impacto sobre el conjunto de la economía.

    \item \textbf{Los encadenamientos productivos amplifican el impacto de los cambios en la demanda.} El aumento del 20\% en la demanda de Agricultura no solo incrementa su propia producción en un 14.4\%, sino que también afecta positivamente a los sectores de Tecnología (+2.0\%) y Servicios (+1.1\%).

    \item \textbf{La comparación de escenarios es útil para la planificación económica.} La evaluación de escenarios conservador, base y optimista permite a los formuladores de políticas anticipar el impacto de diferentes situaciones económicas y tomar decisiones informadas.

    \item \textbf{Los resultados pueden servir como insumo para la formulación de políticas públicas} orientadas al desarrollo económico de La Guajira, identificando sectores estratégicos y evaluando el impacto de programas de estímulo a la demanda.
\end{enumerate}

\clearpage
